
@@M0@@

La
démonstration, immédiate sur la définition fonctorielle des schémas de drapeaux, est laissée au lecteur.

\textbf{Lemme 4.3.} \uline{Soient} @@M1@@ \uline{un anneau commutatif avec unité}, @@M2@@ \uline{un entier} @@M3@@, @@M4@@ \uline{une famille d'éléments de} @@M5@@ ; \uline{soit} @@M6@@ \uline{la} @@M7@@-\uline{algèbre engendrée par des générateurs} @@M8@@, @@M9@@, \uline{soumis aux relations}~:
@@M10@@
\uline{où les} @@M11@@ \uline{désignent les fonctions symétriques élémentaires}. \uline{Soient} @@M12@@ \uline{des entiers} @@M13@@, \uline{de somme} @@M14@@, \uline{et pour} @@M15@@,
@@M16@@
\uline{avec}
@@M17@@
\uline{Posons}~:
@@M18@@
\uline{et soit} @@M19@@ \uline{la sous-algèbre de} @@M20@@ \uline{engendrée par les} @@M21@@. \uline{Alors} :
\begin{enumerate}[label=\textup{\roman*)}]
\item \uline{les éléments de la forme}~:
@@M22@@
\uline{où} @@M23@@, et @@M24@@, \uline{forment une base de} @@M25@@ \uline{considéré comme} @@M26@@-\uline{module}.
\item \uline{les générateurs} @@M27@@ \uline{de la} @@M28@@-\uline{algèbre} @@M29@@ \uline{satisfont aux} @@M30@@ \uline{relations isobares déduites de la relation}~:
@@M31@@
